\documentclass{article}
\usepackage{graphicx} % Required for inserting images
\usepackage{listings}
\usepackage{xcolor}
\usepackage{amsmath}
\usepackage{amssymb}
\usepackage{tikz}
\usetikzlibrary{positioning, arrows.meta}
\usepackage{amsthm}
\usepackage{booktabs}
\newtheorem{example}{Example}
\definecolor{codegreen}{rgb}{0,0.6,0}
\definecolor{codegray}{rgb}{0.5,0.5,0.5}
\definecolor{codepurple}{rgb}{0.58,0,0.82}
\definecolor{backcolour}{rgb}{0.95,0.95,0.92}

\lstdefinestyle{mystyle}{
    backgroundcolor=\color{backcolour},
    commentstyle=\color{codegreen},
    keywordstyle=\color{magenta},
    numberstyle=\tiny\color{codegray},
    stringstyle=\color{codepurple},
    basicstyle=\ttfamily\footnotesize,
    breakatwhitespace=false,
    breaklines=true,
    captionpos=b,
    keepspaces=true,
    numbers=left,
    numbersep=5pt,
    showspaces=false,
    showstringspaces=false,
    showtabs=false,
    tabsize=2
}

\lstset{style=mystyle}
\title{Robust networks}
\author{Racquel Dennison}
\date{March 2026}

\begin{document}

\maketitle

\section{Introduction}

Logistics networks are inherently vulnerable to disruption. When a
transport resource fails mid-operation, whether due to vehicle
breakdown, route closure, or capacity loss this affects the demand serviced within the network as goods already assigned
to failed resources cannot be delivered, and downstream demand goes
unmet. The question of how to plan routes and load assignments that
remain serviceable under such failures is central to supply chain
resilience~\cite{christopher2004building, tomlin2006value}.

We study this question in the context of multi-batching logistics
optimisation, where a network of supply and demand locations is
connected by routing paths, and goods must be transported using
discrete transport resources subject to capacity constraints.
We adopt a two-stage decomposition approach: Stage~1 determines the
flow assignment and load distribution across transport resources,
while Stage~2 handles the bin-packing of items into discrete trips.
The resilience of the overall system is largely determined by
Stage~1, if load is concentrated on a small number of TRs, any
single failure causes disproportionate disruption.

While stage 1 focuses on the distribution of parts within the network, stage 2 focuses on how parts are packed and will be mainly looking at how we can diversify the packed parts within the system.

\section{Related work}


\section{Problem formulation}
\subsection{Problem parameters}

The logistics system is modeled as a multi-commodity flow problem on a structured network. We define the problem as follows:

\subsubsection{Network Topology}

\begin{itemize}
    \item The global logistics network is defined by $N$, where $N = (L, TL)$ :
    \item The set of locations are defined by $L$ , where $l_i \in L$ for $i >0$.
    \item $TL$ denotes the set of transport links. A transport link represents an abstract edge between two locations, $l_i,l_j$. Each link $tl \in TL$ is defined as a triple $tl = (l_i, l_j, R_{tl})$, where:
    \begin{itemize}
        \item $l_i \in L$ is the origin location.
        \item $l_j \in L$ is the destination location.
        \item $R_{tl}$ is the non-empty set of concrete routes available for that specific link. Where each $r \in R_{tl}$ represents a specific route, $r = (tr,d_r)$, indicates the use of transport resource $tr$ over distance $d_r$.
    \end{itemize}
\end{itemize}

\subsubsection{Transportation}
\begin{itemize}
    \item $TR$ denotes the set of available transport resources. Each $tr \in TR$ is defined by a tuple $(cost, speed, co2emissions, capacity)$:
    \begin{itemize}
        \item $cost$ : Operating cost per unit distance.
        \item $speed$ : Average transport speed.
        \item $co2emissions$ : $CO_2$ emissions per unit distance.
        \item $capacity$ : Maximum capacity of the resource.
    \end{itemize}
\end{itemize}

\subsubsection{Products}
\begin{itemize}
    \item $P$ denotes the set of products being transported in the network $p$. Each product $p \in P$ is defined by $(weight, value, TR_p)$:
    \begin{itemize}
        \item $weight$ : Unit size/volume, in this case here, we are only considering a one dimensional size.
        \item $value$ : The unit monetary value.
        \item $TR_p \subseteq TR$ : Set of valid transport resources for product $p$ to be transported on.
    \end{itemize}
\end{itemize}

\subsubsection{Packings and Decision Variables}
We represent the possible resource packings by the following sets.
\begin{itemize}
    \item $\mathcal{B}_r$ : The set of valid transport resource packings for route $r$.
    \item $B = \{ (p, q_{pb}) : p \in P \} \in \mathcal{B}_r$ : A packing configuration where $quanity_{pb}$ is the amount of product $p$ packed into batch $B$, satisfying $\sum_{p \in P} quanity_{pb} *  weight \leq capacity$.
    \item $\lambda_B \in \mathbb{N}$ : Decision variable representing the frequency of usage for packing $B$.
\end{itemize}

A valid solution to the multi batching routing problem consists of the enumeration of the set packages assigned to a route denoted as $\mathcal{B}_r$. Furthermore, the most optimal valid solution aims to reduce the overall cost of the network as much as possible.



\begin{example}[Running Example]\label{ex:running}
Consider the following logistics network with the locations $\{l_1,l_2,l_3,l_4\} $which has a single
supply node $l_1$, a single demand node $l_4$, and two intermediate
hub nodes $l_2$ and $l_3$. Two transport resources are available:
$tr_1$  = (cost 40 per trip, capacity 8) and $tr_2$ =  (cost 50 per
trip, capacity 8). Two parts need to be transported from $l_1$ to $l_4$, with each part having a size of 2.

\begin{figure}[h]
\centering
\begin{tikzpicture}[
    node/.style   = {circle, draw, thick, minimum size=1cm, font=\small},
    supply/.style = {node, fill=green!15},
    demand/.style = {node, fill=red!15},
    hub/.style    = {node, fill=gray!10},
    lbl/.style    = {font=\scriptsize, fill=white, inner sep=2pt}
]

% Nodes
\node[supply] (l1) at (0, 0)    {$l_1$};
\node[hub]    (l2) at (3, 1.5)  {$l_2$};
\node[hub]    (l3) at (3,-1.5)  {$l_3$};
\node[demand] (l4) at (6, 0)    {$l_4$};

% l1 -> l2 (Split TRs)
\draw[-latex, thick, dotted] (l1) to[bend left=15]
    node[lbl, above sloped] {$t_1$: 40} (l2);
\draw[-latex, thick] (l1) to[bend right=15]
    node[lbl, below sloped] {$t_2$: 50} (l2);

% l2 -> l4 (Split TRs)
\draw[-latex, thick, dotted] (l2) to[bend left=15]
    node[lbl, above sloped] {$t_1$: 40} (l4);
\draw[-latex, thick] (l2) to[bend right=15]
    node[lbl, below sloped] {$t_2$: 50} (l4);

% l1 -> l3 (Split TRs)
\draw[-latex, thick, dotted] (l1) to[bend left=15]
    node[lbl, above sloped] {$tr_1$: 60} (l3);
\draw[-latex, thick] (l1) to[bend right=15]
    node[lbl, below sloped] {$tr_2$: 70} (l3);

% l3 -> l4 (Split TRs)
\draw[-latex, thick, dotted] (l3) to[bend left=15]
    node[lbl, above sloped] {$tr_1$: 60} (l4);
\draw[-latex, thick] (l3) to[bend right=15]
    node[lbl, below sloped] {$tr_2$: 70} (l4);

% Supply / demand annotations
% Adjusted position slightly so it doesn't bump into the new bent lines
\node[left=0.2cm of l1, font=\scriptsize]
    {supply: $p_1{=}8,\, p_2{=}8$};
\node[right=0.2cm of l4, font=\scriptsize]
    {demand: $p_1{=}8,\, p_2{=}8$};

\end{tikzpicture}
\caption{Network flow diagram with separated transport resources and colored node classifications.}
\end{figure}

\end{example}
\subsection{Stage 1: Flow Assignment without soft constraints}

Stage~1 determines, for each active route $(f, t, tr)$ and each part
$p \in P$, the \emph{load} $\ell(f, t, tr, p) \in \mathbb{Z}_{\geq 0}$ represents the number of units of $p$ transported along that route via $tr$.
It also selects a \emph{trip frequency} $\phi(f, t, tr) \in \mathbb{Z}_{>0}$
for each active route, representing how many trips $tr$ makes on arc
$(f, t)$.

The decision variables are subject to three classes of constraint.

\paragraph{Flow conservation.} For each part $p$ and location $l$, the
net flow must satisfy the supply/demand requirement:
\[
    \sum_{tr,\, t:\, (l,t,tr) \in R} \ell(l,t,tr,p)
    \;-\;
    \sum_{tr,\, f:\, (f,l,tr) \in R} \ell(f,l,tr,p)
    \;=\; \delta(p, l)
    \quad \forall\, p \in P,\, l \in L.
\]

\paragraph{Capacity.} The total size of the parts assigned to some transport resource must not exceed the total capacity of this transport resource:
\[
    \sum_{p \in P} s(p)\cdot \ell(f, t, tr, p)
    \;\leq\; \tfrac{4}{5}\cdot \kappa(tr) \cdot \phi(f, t, tr)
    \quad \forall\, (f, t, tr) \in R.
\]

\paragraph{Objective.} Stage~1 minimises total transport cost:
\[
    \min \sum_{(f,t,tr) \in R} c(f, t, tr) \cdot \phi(f, t, tr).
\]

In \textsc{Clingcon}, we model these three constraints above with the use of theory terms. This is outlined in Listing \ref{theory_stage1_constraints}.

\begin{lstlisting}[language=Prolog, caption={Stage 1 core constraints}, label={theory_stage1_constraints}]
% Flow domain
&dom { 0..Max } = load(F, T, TR, P) :-
    possibleFlow(F, T, TR, P), totalSupply(P, Max).

% Flow conservation
&sum { flow_in(P,L); -flow_out(P,L) } = Net :-
    requiredNet(P, L, Net).

% 80% capacity headroom
&sum { S * load(F, T, TR, P) :
        possibleFlow(F, T, TR, P), partSizeSmall(P, S) }
    <= TotalCap :-
    routeFreq(F, T, TR, _, TotalCap).

% Cost minimisation
#minimize { V, F, T, TR : _transCost(F, T, TR, V) }.
\end{lstlisting}
Consider the example in \ref{ex:running}, we can first represent the facts of the network as follows:
\begin{lstlisting}[language=Prolog, caption={Stage 1 core constraints}, label={example_instances}]

location(l1).
location(l2).
location(l3).
location(l4).


transportResource(tr1).
transportResource(tr2).

transportCapacity(tr1, 8).
transportCapacity(tr2, 8).


part(p1).
part(p2).

partSize(p1, 2).
partSize(p2, 2).

partVal(p1, 500).
partVal(p2, 400).

partTR(p1, tr1).
partTR(p1, tr2).
partTR(p2, tr1).
partTR(p2, tr2).

demandOffer(p1, l1,  8).
demandOffer(p2, l1,  8).
demandOffer(p1, l4, -8).
demandOffer(p2, l4, -8).


route(l1, l2, tr1, 1, 40).
route(l1, l2, tr2, 1, 50).
route(l2, l4, tr1, 1, 40).
route(l2, l4, tr2, 1, 50).
route(l1, l3, tr1, 2, 60).
route(l1, l3, tr2, 2, 70).
route(l3, l4, tr1, 2, 60).
route(l3, l4, tr2, 2, 70).

\end{lstlisting}

If we computed stage one against this example above, we would obtain the results shown in Table \ref{tab:routing-assignments}. If we look at the results, we see that none of the routing goes through $l3$. If for example, the transport resource line for $t_1$ from $l_1 \rightarrow l_2$ were to fail, a huge majority of the network would not be able to service the demand. We can actually caculate how much demand would be serviced in this instance.
\begin{table}[h]
\centering
\begin{tabular}{|l|l|l|c|}

\hline
\textbf{Route} & \textbf{Transport Resource} & \textbf{Load} & \textbf{Frequency} \\
\hline
$l_1 \to l_2$ & $t1$ & $p_1=7,\; p_2=5$ & 4 \\
\hline
$l_1 \to l_2$ & $t2$ & $p_1=1,\; p_2=3$ & 2 \\
\hline
$l_2 \to l_4$ & $t1$ & $p_1=6,\; p_2=6$ & 4 \\
\hline
$l_2 \to l_4$ & $t2$ & $p_1=2,\; p_2=2$ & 2 \\
\hline

\end{tabular}
\caption{Optimal routing assignments (Total Cost: 520). Routes with zero load and frequency are omitted.}
\label{tab:routing-assignments}
\end{table}
\subsubsection{Transport Resilience Metric $R_{TR}$}
\label{sec:rtr}

Stage~1 produces a least-cost load assignment, but as observed in
the running example, the resulting plan may concentrate load on a
small number of transport resources. In order to push the solver towards distributing loads over transport resources, we can define a $\texttt{dominantTR}$ to be the transport resource which is carrying more than what can be allocated as load on a given arc. We can list it as follows:
\begin{lstlisting}
    dominantTR(F, T, P) :-
    transportResourcesAvailable(F, T, P, N), N > 1,
    &sum { N*maxTRLoad(F, T, P); -arcLoad(F, T, P) } > 0.
\end{lstlisting}

In order to evaluate the effect of applying this soft constraint, we will need some sort of metric to measure the effects. For a given Stage~1 solution, we define the \emph{transport
resilience} $R_{TR}$ as the worst-case fraction of total demand
that can still be served under any single transport resource
failure, assuming the operator responds optimally by rerouting
through all surviving arcs.

Formally, given a set of active arcs $\mathcal{A}$ and a disruption
$(f^*, t^*, tr^*) \in \mathcal{A}$, we construct a flow network
$\mathcal{N}_p^{(f^*,t^*,tr^*)}$ for each part $p \in P$ as
follows:

\begin{itemize}
    \item A super-source node $S$ with arc $S \to l$ of capacity $\delta(p,l)$ for each supply location $l$.
    \item A super-sink node $T$ with arc $l \to T$ of capacity $-\delta(p,l)$ for each demand location $l$.
    \item For each active arc $(f,t,tr) \in \mathcal{A}$, an internal arc $f \to t$ with capacity:
    \[
        \mathrm{cap}_D(f,t,tr,p) =
        \begin{cases}
            0 & \text{if } (f,t,tr) = (f^*,t^*,tr^*) \\
            load(f,t,tr,p) & \text{otherwise}
        \end{cases}
    \]
\end{itemize}

The \emph{coverage} of disruption $(f^*,t^*,tr^*)$ for part $p$
is the fraction of demand that can be satisfied via max-flow in the
surviving network:

\[
    \mathrm{cov}(f^*,t^*,tr^*,p) =
    \frac{\mathrm{maxflow}\!\left(
        \mathcal{N}_p^{(f^*,t^*,tr^*)},\, S,\, T
    \right)}{\mathrm{totalDemand}(p)}
\]

Since parts are co-required at demand nodes, the effective coverage
is limited by the scarcest part. We therefore take the minimum over
all parts. The overall $R_{TR}$ is the minimum over all possible
disruptions:

\begin{equation}\label{eq:rtr}
    R_{TR} = \min_{(f,t,tr) \in \mathcal{A}}\;
             \min_{p \in P}\;
             \mathrm{cov}(f,t,tr,p)
\end{equation}

$R_{TR} \in [0,1]$, where $R_{TR} = 1$ means every single-TR
failure can be fully absorbed, and $R_{TR} = 0$ means at least one
failure leaves demand entirely unserviceable.

\textbf{Worked Example}

Returning to Example~\ref{ex:running} with the Stage~1 solution
from Table~\ref{tab:routing-assignments}, we compute $R_{TR}$ by
evaluating every active arc-TR disruption.

\paragraph{Disruption: lose $tr_1$ on $l_1 \to l_2$.}

$tr_1$ was the dominant carrier on this arc: it carried
$p_1 = 7$ and $p_2 = 5$ of the 8 units demanded. After the
disruption, only $tr_2$'s load survives on this arc ($p_1=1$,
$p_2=3$). We build $\mathcal{N}_{p_1}^{(l_1,l_2,tr_1)}$:

\[
\begin{array}{ll}
S \to l_1: & \mathrm{cap} = 8 \\
l_1 \to l_2 \text{ via } tr_1: & \mathrm{cap} = 0
    \quad \text{(disrupted)}\\
l_1 \to l_2 \text{ via } tr_2: & \mathrm{cap} = 1 \\
l_2 \to l_4 \text{ via } tr_1: & \mathrm{cap} = 6 \\
l_2 \to l_4 \text{ via } tr_2: & \mathrm{cap} = 2 \\
l_4 \to T: & \mathrm{cap} = 8
\end{array}
\]

The only path from $S$ to $T$ is $S \to l_1 \to l_2
(tr_2) \to l_4 \to T$. The bottleneck is
$\min(8, 1, 8, 8) = 1$, so $\mathrm{maxflow} = 1$.

\[
    \mathrm{cov}(l_1, l_2, tr_1, p_1) = \frac{1}{8} = 0.125
\]

Repeating for $p_2$ (surviving $tr_2$ load = 3):

\[
    \mathrm{cov}(l_1, l_2, tr_1, p_2) = \frac{3}{8} = 0.375
\]

The joint coverage for this disruption is
$\min(0.125, 0.375) = 0.125$.

\paragraph{All disruptions.}

Table~\ref{tab:rtr-disruptions} reports the coverage for every
active arc-TR disruption. The worst case is losing $tr_1$ on
$l_1 \to l_2$, giving:

\[
    R_{TR} = 0.125
\]

\begin{table}[h]
\centering
\begin{tabular}{|l|l|c|r|r|}
\hline
\textbf{Disruption} & \textbf{Part} &
\textbf{Max-flow} & \textbf{Demand} & \textbf{Coverage} \\
\hline
$(l_1,l_2,tr_1)$ & $p_1$ & 1 & 8 & 0.125 \\
$(l_1,l_2,tr_1)$ & $p_2$ & 3 & 8 & 0.375 \\
$(l_1,l_2,tr_2)$ & $p_1$ & 7 & 8 & 0.875 \\
$(l_1,l_2,tr_2)$ & $p_2$ & 5 & 8 & 0.625 \\
$(l_2,l_4,tr_1)$ & $p_1$ & 1 & 8 & 0.125 \\
$(l_2,l_4,tr_1)$ & $p_2$ & 3 & 8 & 0.375 \\
$(l_2,l_4,tr_2)$ & $p_1$ & 7 & 8 & 0.875 \\
$(l_2,l_4,tr_2)$ & $p_2$ & 5 & 8 & 0.625 \\
\hline
\multicolumn{4}{l}{$R_{TR}$ (minimum over all disruptions
and parts)} & \textbf{0.125} \\
\hline
\end{tabular}
\caption{Coverage values for all active arc-TR disruptions
in the running example. The worst case arises when $tr_1$
fails on either $l_1 \to l_2$ or $l_2 \to l_4$, leaving
only 1 unit of $p_1$ serviceable $12.5\%$ of demand.}
\label{tab:rtr-disruptions}
\end{table}

This low value reflects the concentration of load on $tr_1$:
at $w=0$, Stage~1 assigned 7 of 8 units of $p_1$ to $tr_1$
on $l_1 \to l_2$, leaving $tr_2$ with only 1 unit. Any
failure of $tr_1$ is therefore catastrophic for $p_1$
coverage.

Lets say that we compute this same example again but we introduce the $\texttt{dominantTR}$ predicate. We also have the following soft constraint in order to minimise the total $\texttt{dominantTR}$ arcs.
\begin{lstlisting}
    :~ dominantTR(F, T, P). [weight@1, div, F, T, P]
\end{lstlisting}

If we set the weight to $50$, Stage~1 produces the solution shown
in Table~\ref{tab:routing-assignments-w50}. The solver now splits
load evenly between $tr_1$ and $tr_2$ on every arc: each carries
exactly 4 units of $p_1$ and 4 units of $p_2$ on both
$l_1 \to l_2$ and $l_2 \to l_4$. No \textsc{dominantTR} atom
fires. The total cost increases modestly from 520 to 540, reflecting
the use of $tr_2$ which costs 50 per trip versus 40 for $tr_1$.

\begin{table}[h]
\centering
\begin{tabular}{|l|l|c|c|}
\hline
\textbf{Route} & \textbf{TR} & \textbf{Load} & \textbf{Freq} \\
\hline
$l_1 \to l_2$ & $tr_1$ & $p_1=4,\; p_2=4$ & 3 \\
$l_1 \to l_2$ & $tr_2$ & $p_1=4,\; p_2=4$ & 3 \\
$l_2 \to l_4$ & $tr_1$ & $p_1=4,\; p_2=4$ & 3 \\
$l_2 \to l_4$ & $tr_2$ & $p_1=4,\; p_2=4$ & 3 \\
\hline
\end{tabular}
\caption{Optimal routing assignments at $w=50$ (Total Cost: 540).
Load is perfectly balanced across $tr_1$ and $tr_2$ on every arc.
Routes with zero load are omitted.}
\label{tab:routing-assignments-w50}
\end{table}

We can immediately observe the effect on $R_{TR}$. If $tr_1$ now
fails on $l_1 \to l_2$, the surviving $tr_2$ load is $p_1=4$ and
$p_2=4$ exactly half of total demand. The max-flow network
for $p_1$ after losing $(l_1, l_2, tr_1)$ yields:

\[
    \mathrm{cov}(l_1, l_2, tr_1, p_1) = \frac{4}{8} = 0.5
\]

By symmetry, every arc-TR disruption produces the same coverage
of $0.5$ across both parts. Therefore:

\[
    R_{TR} = 0.5
\]

\section{Experiments}

To test the relationship between the cost of the network and the diversification constraint, we ran a series of experiments. The experiments were based on a larger network encoding with 6 locations, 2 parts and 2 transport resources. We varied the weight $w$ of the \texttt{dominantTR} soft constraint across 18 values from $w=0$ (pure cost minimisation) to $w=1{,}000$ (strong diversification pressure). Our hypotheses were: (i)~increasing $w$ produces a tradeoff with transport cost, and (ii)~as $w$ increases, $R_{TR}$ improves because no single transport resource dominates the transportation of goods.

\subsection{Experimental Setup}

The test instance consists of 6 locations ($l_1$--$l_6$), 2 transport resources ($tr_1$ with capacity~10 and cost~56, $tr_2$ with capacity~15 and cost~42), 2 parts ($p_1$ with size~4, $p_2$ with size~3), and 17 directed routes. Total supply is 30 units of $p_1$ and 40 units of $p_2$, with demand distributed across locations $l_2$, $l_5$, and $l_6$.

Each configuration was solved using Clingcon with a time limit of 600 seconds. For each solution, we extracted:
\begin{itemize}
    \item The \emph{transport cost}: the pure routing cost, computed as the nominal objective value minus the penalty contribution ($w \times |\texttt{dominantTR}|$).
    \item The \emph{dominant count}: the number of $(arc, part)$ pairs where one TR carries more than its fair share.
    \item $R_{TR}$: computed post-hoc via the Edmonds--Karp max-flow algorithm over all single-TR disruption scenarios.
\end{itemize}

The weight was swept over $w \in \{0, 10, 15, 20, 25, 30, 40, 50, 75, 100, 125, 150, 200, 300, 500, 600, 750, 1000\}$.

\subsection{Results}

Table~\ref{tab:pareto-results} presents the results for selected weight values that illustrate the key transitions in the cost--resilience tradeoff.

\begin{table}[h]
\centering
\begin{tabular}{|r|r|r|c|l|c|}
\hline
$w$ & \textbf{Transport Cost} & \textbf{Dominant Count} & $R_{TR}$ & \textbf{Worst Arc} & \textbf{Optimal?} \\
\hline
0    & 5{,}040 & 8 & 0.500 & $(l_1,l_3,tr_2)$ & No \\
\hline
10   & 4{,}466 & 7 & 0.600 & $(l_1,l_3,tr_2)$ & No \\
\hline
50   & 4{,}466 & 7 & 0.600 & $(l_1,l_2,tr_2)$ & No \\
\hline
100  & 4{,}550 & 6 & 0.500 & $(l_1,l_3,tr_2)$ & No \\
\hline
150  & 5{,}040 & 4 & 0.500 & $(l_1,l_3,tr_2)$ & No \\
\hline
300  & 5{,}236 & 3 & 0.600 & $(l_1,l_3,tr_2)$ & No \\
\hline
500  & 6{,}216 & 0 & 0.667 & $(l_2,l_6,tr_1)$ & No \\
\hline
750  & 6{,}216 & 0 & 0.667 & $(l_2,l_6,tr_1)$ & Yes (456\,s) \\
\hline
1000 & 6{,}216 & 0 & 0.667 & $(l_2,l_6,tr_1)$ & Yes (87\,s) \\
\hline
\end{tabular}
\caption{Selected results from the \texttt{div\_weight} sweep on the 6-location instance. Transport cost is the pure routing cost (excluding penalty). Dominant count is the number of \texttt{dominantTR} atoms in the solution. $R_{TR}$ is the worst-case coverage under any single-TR failure, computed via max-flow.}
\label{tab:pareto-results}
\end{table}


\subsection{Analysis}

\subsubsection{Dominant count decreases monotonically with weight}

As $w$ increases, the solver progressively eliminates dominant transport resource assignments. The dominant count drops from 8 at $w=0$ through intermediate values (7 at $w=10$--$75$, then 6, 4, 3) down to 0 at $w \geq 500$. At $w \geq 500$, no single TR carries more than its fair share on any arc---the load is perfectly balanced across transport resources. This confirms that the \texttt{dominantTR} penalty is effective at driving the solver toward load-balanced solutions.

\subsubsection{The ``free lunch'' region: $w = 10$--$75$}

A surprising finding is that moderate diversification pressure actually \emph{reduces} transport cost relative to the unconstrained baseline. At $w=0$ the transport cost is 5{,}040, but at $w=10$--$75$ it drops to 4{,}466---an 11.4\% reduction. This occurs because the penalty steers the solver away from locally attractive but globally suboptimal routing decisions within the 600-second time limit. The diversification constraint effectively acts as a search guidance heuristic: by penalising concentrated solutions, it redirects the solver toward regions of the search space that happen to contain lower-cost alternatives that the pure cost objective did not find in time.

In this region, $R_{TR}$ simultaneously improves from 0.500 to 0.600, meaning the network can guarantee at least 60\% demand coverage under any single-TR failure---up from 50\%. This represents a Pareto improvement: strictly better on both cost and resilience.

\subsubsection{The cost--resilience tradeoff emerges at higher weights}

Beyond $w=75$, the expected tradeoff materialises. Transport cost rises from 4{,}466 through 4{,}550 ($w=100$), 5{,}040 ($w=150$), and 5{,}236 ($w=300$) to 6{,}216 ($w \geq 500$). The maximum resilience of $R_{TR} = 0.667$ is achieved at $w \geq 500$, where the dominant count reaches zero. This represents the theoretical upper bound for a two-TR system: with perfectly balanced loads, losing either TR leaves exactly $1 - 1/n_{\text{arcs}}$ of the capacity surviving, where the bottleneck arc determines the network-wide minimum.

The cost premium for full diversification relative to the best transport cost found ($w=10$) is:
\[
    \frac{6{,}216 - 4{,}466}{4{,}466} = 39.2\%
\]

\subsubsection{$R_{TR}$ is non-monotone with weight}

While the dominant count decreases monotonically, $R_{TR}$ does not increase monotonically. At $w=100$--$150$, $R_{TR}$ dips back to 0.500 despite having fewer dominant arcs (6 and 4 respectively) than the $w=10$ solution (7 dominant arcs). This reveals that the \emph{number} of dominant arcs is an imperfect proxy for resilience: it is possible to have fewer imbalanced arcs overall while still concentrating load heavily on the particular arc that constitutes the network bottleneck. $R_{TR}$ captures this because it evaluates the worst-case disruption across the entire network topology via max-flow, not just a local arc-level statistic.

\subsubsection{Higher weights improve solver convergence}

An additional observation is that the solver reaches optimality significantly faster at high weights. At $w \leq 600$, no configuration converges within the 600-second limit. At $w=750$ the solver proves optimality in 456 seconds, and at $w=1{,}000$ in just 87 seconds. This suggests that the diversification penalty simplifies the search landscape: by strongly penalising the many concentrated solutions, it effectively prunes large regions of the search space, allowing the solver to converge to the (unique) fully balanced configuration more quickly.

\subsection{Pareto Front}

Figure~\ref{fig:pareto-front} illustrates the Pareto front in the transport cost versus dominant count space. The non-dominated solutions are:

\begin{table}[h]
\centering
\begin{tabular}{|r|r|r|c|}
\hline
$w$ & \textbf{Transport Cost} & \textbf{Dominant Count} & $R_{TR}$ \\
\hline
10  & 4{,}466 & 7 & 0.600 \\
\hline
100 & 4{,}550 & 6 & 0.500 \\
\hline
150 & 5{,}040 & 4 & 0.500 \\
\hline
300 & 5{,}236 & 3 & 0.600 \\
\hline
500 & 6{,}216 & 0 & 0.667 \\
\hline
\end{tabular}
\caption{Pareto-optimal configurations minimising both transport cost and dominant count. Each row is non-dominated: no other configuration achieves both lower cost and fewer dominant arcs.}
\label{tab:pareto-front}
\end{table}

The Pareto curve reveals a non-linear cost structure. The first step---reducing from 7 to 6 dominant arcs---costs only 84 units ($+1.9\%$). However, eliminating the final 3 dominant arcs (from 3 to 0) costs 980 units ($+18.7\%$). This suggests diminishing marginal returns: the last few dominant arcs are the most expensive to eliminate because they correspond to arcs where the asymmetry between $tr_1$ and $tr_2$ (different capacities and costs) makes balanced allocation inherently costly.

\subsection{Discussion}

The results confirm both hypotheses with important nuances:

\paragraph{Hypothesis (i): Cost--resilience tradeoff.} Confirmed for $w > 75$. However, the ``free lunch'' region ($w = 10$--$75$) shows that moderate diversification pressure can simultaneously reduce cost and improve resilience. This occurs because the penalty provides useful search guidance to the solver within the time limit.

\paragraph{Hypothesis (ii): $R_{TR}$ improves with weight.} Confirmed in the aggregate---$R_{TR}$ improves from 0.500 ($w=0$) to 0.667 ($w \geq 500$)---but the path is non-monotone. The dominant count is a useful proxy but not a perfect predictor of $R_{TR}$. The network-level max-flow computation reveals that the \emph{topology} of load concentration matters, not just the count of imbalanced arcs.

\paragraph{Practical implication.} For this instance, a practitioner could set $w \approx 10$--$50$ and obtain both lower cost and better resilience than the unconstrained baseline. If maximum resilience is required, $w \geq 500$ achieves $R_{TR} = 0.667$ at a 39\% cost premium. The choice of $w$ thus depends on the operator's risk tolerance and the relative cost of service disruption versus additional transport expenditure.


\end{document}
